\documentclass[11pt]{article}

\usepackage[margin=1in]{geometry}
\usepackage{amsmath,amssymb,amsthm,mathtools}
\usepackage{booktabs}
\usepackage{enumitem}
\usepackage[colorlinks=true,linkcolor=blue,urlcolor=blue]{hyperref}
\usepackage{microtype}

\newtheorem{theorem}{Theorem}[section]
\newtheorem{proposition}[theorem]{Proposition}
\newtheorem{lemma}[theorem]{Lemma}
\newtheorem{corollary}[theorem]{Corollary}
\newtheorem{conjecture}[theorem]{Conjecture}
\theoremstyle{definition}
\newtheorem{definition}[theorem]{Definition}
\newtheorem{remark}[theorem]{Remark}

\newcommand{\PP}{\mathbb P}
\newcommand{\E}{\mathbb E}
\newcommand{\con}{\leftrightarrow}
\newcommand{\ncon}{\nleftrightarrow}
\newcommand{\1}{\mathbf 1}
\newcommand{\argmin}{\operatorname*{arg\,min}}

\title{Kozma--Nitzan Question 8:\\
exact statement, boundary counterexample, and a formalization-ready core}
\author{}
\date{September 2, 2026}

\begin{document}
\maketitle

\begin{abstract}
The printed Question 8 in Kozma--Nitzan minimizes
\(\PP(a\con b,\,o\ncon A)\), not
\(\PP(o\con A,\,a\con b)\).  This distinction is essential.  Under the
usual reading that the desired inequality must hold for every chosen
minimizer, the printed question is false when edge probabilities in
\([0,1]\) are allowed: an exact four-vertex example is given below.  The
unique-minimizer and positive-avoidance versions are not refuted by that
example and remain open here.

This note gives a statement-faithful formulation, the complete
counterexample, several proved reductions, and an atomized proof interface
for Lean.  In particular, the global unique-minimizer version is equivalent
to the version for every minimizer under the hypothesis
\(\PP(o\ncon A)>0\).  The reduction uses a support-graph lemma for a
zero score and a pendant-relay tie breaker for a positive score.  No open
claim is used as a proved lemma.
\end{abstract}

\section{Model and exact notation}

Let \(G=(V,E)\) be a finite graph.  Every edge \(e\in E\) is independently
open with probability \(p_e\in[0,1]\).  Connectivity in the random open
subgraph is denoted by \(x\con y\), and connectivity to a nonempty finite
set \(A\subseteq V\) by \(x\con A\).  Connectivity is reflexive.  Fix
vertices \(o,b\in V\), and put
\begin{equation}\label{eq:U}
 U:=\{o\con A\},\qquad D:=U^{\mathrm c}=\{o\ncon A\}.
\end{equation}
For every \(x\in A\), define
\begin{equation}\label{eq:three-quantities}
 L_x:=\PP(U,\,x\con b),\qquad
 \rho_x:=\PP(D,\,x\con b),\qquad
 R:=\PP(U,\,o\con b).
\end{equation}
The elementary partition \(U\sqcup D\) gives
\begin{equation}\label{eq:total-split}
 \PP(x\con b)=L_x+\rho_x\qquad(x\in A).
\end{equation}

With this notation, display (41) in the paper is \(L_a\le R\), while the
selection rule in its Question 8 is
\begin{equation}\label{eq:q8-selection}
 a\in\argmin_{x\in A}\rho_x
 =\argmin_{x\in A}\PP(x\con b,\,o\ncon A).
\end{equation}
The slash in \(o\ncon A\) is visible in the rendered original PDF.  Text
extraction from that PDF can display the glyph incorrectly as ``\(o=A\)'',
which is the source of the earlier transcription error.

\section{The formulations that must not be conflated}

The phrase ``let \(a\) be the point minimising'' is ambiguous when there
is a tie.  The following statements make the possible quantifiers explicit.

\begin{definition}[Every-minimizer form]\label{def:every}
For every finite graph \(G=(V,E)\), every weight vector
\((p_e)_{e\in E}\in[0,1]^E\), every nonempty \(A\subseteq V\), every
\(o,b\in V\), and every \(a\in A\),
\begin{equation}\label{eq:every}
 a\in A,\quad \rho_a\le \rho_x\ \text{for every }x\in A
 \quad\Longrightarrow\quad L_a\le R.
\end{equation}
\end{definition}

\begin{definition}[Strict-minimizer form]\label{def:strict}
For every finite graph \(G=(V,E)\), every weight vector
\((p_e)_{e\in E}\in[0,1]^E\), every nonempty \(A\subseteq V\), every
\(o,b\in V\), and every \(a\in A\),
\begin{equation}\label{eq:strict}
 a\in A,\quad \rho_a<\rho_x\ \text{for every }x\in A\setminus\{a\}
 \quad\Longrightarrow\quad L_a\le R.
\end{equation}
\end{definition}

\begin{definition}[Positive-avoidance form]\label{def:positive}
With the same closed quantifiers as in Definition~\ref{def:every}, the
implication \eqref{eq:every} is required under the additional hypothesis
\(\PP(D)>0\).
\end{definition}

\begin{definition}[Existential form]\label{def:exists}
For every finite graph \(G=(V,E)\), every weight vector in \([0,1]^E\),
every nonempty \(A\subseteq V\), and every \(o,b\in V\), there is at least
one \(a\in\argmin_{x\in A}\rho_x\) for which \(L_a\le R\).
\end{definition}

\begin{center}
\begin{tabular}{@{}ll@{}}
\toprule
Formulation & Status established in this note \\
\midrule
Every minimizer, probabilities in \([0,1]\) & False; Theorem~\ref{thm:counterexample} \\
Strict or unique minimizer & Open \\
Every minimizer with \(\PP(D)>0\) & Open; equivalent to the strict form \\
Some minimizing relay is good & Open in general \\
Minimizer of \(L_x\), rather than \(\rho_x\) & True by strong Conjecture 2 \\
\bottomrule
\end{tabular}
\end{center}

Thus the unqualified assertion ``Question 8 is proved'' is not justified.
The every-minimizer form is refuted, and the other natural readings still
require a new argument.

\section{Exact four-vertex calculation}

The counterexample is most transparent as a symbolic three-parameter
calculation.

\begin{lemma}[Weighted four-vertex path]\label{lem:path}
Let
\[
 V=\{o,a_1,a_2,b\},\qquad A=\{a_1,a_2\},
\]
let \(s,q,t\in[0,1]\), and consider the path
\[
 o\;\text{---}\;a_1\;\text{---}\;a_2\;\text{---}\;b.
\]
Give its three consecutive edges probabilities \(s,q,t\), respectively,
and give every other edge weight zero.  Then
\begin{align}
 \PP(U)&=s,                         & R&=sqt,\label{eq:path1}\\
 L_{a_1}&=sqt,                      & L_{a_2}&=st,\label{eq:path2}\\
 \rho_{a_1}&=(1-s)qt,              & \rho_{a_2}&=(1-s)t.\label{eq:path3}
\end{align}
\end{lemma}

\begin{proof}
Write \(X,Y,Z\) for the indicators that the three consecutive edges are
open.  The vertex \(o\) has only the first edge available, and \(a_1\in A\),
so \(U=\{X=1\}\).  The event \(\{o\con b\}\) is
\(\{X=Y=Z=1\}\).  The connection \(a_1\con b\) requires
\(Y=Z=1\), whereas \(a_2\con b\) requires only \(Z=1\).
Independence now gives all six identities.
\end{proof}

\begin{theorem}[The every-minimizer form is false]\label{thm:counterexample}
Definition~\ref{def:every} is false in the percolation model of the paper.
The failure occurs with four distinct vertices and with \(o,b\notin A\).
\end{theorem}

\begin{proof}
In Lemma~\ref{lem:path}, set
\[
 s=1,\qquad q=\frac12,\qquad t=1.
\]
These weights are admissible because the model permits probabilities in
\([0,1]\).  Equation \eqref{eq:path3} gives
\(\rho_{a_1}=\rho_{a_2}=0\), so both relays are minimizers in
\eqref{eq:q8-selection}.  If the selected minimizer is \(a_2\), however,
\[
 L_{a_2}=1>\frac12=R
\]
by \eqref{eq:path1}--\eqref{eq:path2}.  This is precisely a failure of
display (41).
\end{proof}

\begin{remark}[Why the obvious perturbation is not a strict counterexample]
Keep \(0<s<1\), \(0\le q<1\), and \(t>0\).  Then
\(\rho_{a_1}<\rho_{a_2}\), so \(a_1\) is the unique Question 8 minimizer;
but \(L_{a_1}=R=sqt\).  The displayed boundary counterexample therefore
does not refute Definition~\ref{def:strict}.
\end{remark}

\section{Diagnosis of the earlier argument}

The already proved strong form of Conjecture 2 is
\begin{equation}\label{eq:strongC2}
 R\ge \min_{x\in A}L_x.
\end{equation}
Consequently, a relay minimizing \(L_x=\PP(o\con A,x\con b)\) satisfies
display (41).  The earlier manuscript used exactly this observation, but
it transcribed the selection rule in Question 8 as minimization of \(L_x\).
The printed selection rule is minimization of \(\rho_x\).

Identity \eqref{eq:total-split} says
\[
 \rho_x=\PP(x\con b)-L_x.
\]
It supplies no equivalence between the two rankings.  In the
counterexample, both \(\rho\)-scores vanish while
\(L_{a_1}=1/2\) and \(L_{a_2}=1\).  Hence the strong Conjecture 2 argument
proves a useful neighboring statement, but it proves neither the strict
nor the every-minimizer form of the printed Question 8.

\section{Elementary cases that are settled}

\begin{proposition}[One relay]\label{prop:singleton}
If \(A=\{a\}\), then \(L_a=R\).
\end{proposition}

\begin{proof}
On \(U=\{o\con a\}\), the open clusters of \(o\) and \(a\) coincide.
Thus \(\1_{\{a\con b\}}=\1_{\{o\con b\}}\) on \(U\).
\end{proof}

\begin{proposition}[The degenerate locus]\label{prop:degenerate}
If \(\PP(D)=0\), then every \(x\in A\) minimizes \(\rho_x\).  Moreover,
\eqref{eq:strongC2} guarantees at least one minimizing relay \(x\) with
\(L_x\le R\), but Theorem~\ref{thm:counterexample} shows that another
minimizing relay can fail the inequality.
\end{proposition}

\begin{proof}
The hypothesis makes \(\rho_x=\PP(D,x\con b)=0\) for every \(x\in A\).
The existence assertion is \eqref{eq:strongC2}; the failure of the universal
assertion was computed above.
\end{proof}

\begin{proposition}[Overlap with Question 7]\label{prop:q7}
Suppose \(a\) is a Question 8 minimizer and also satisfies
\[
 \PP(a\con b)\le \PP(x\con b)\qquad(x\in A).
\]
Then \(L_a\le R\).
\end{proposition}

\begin{proof}
This is exactly the already machine-checked affirmative answer to
Question 7, applied to the specified unconditional minimizer.
\end{proof}

\begin{proposition}[Observer as target]\label{prop:b=o}
If \(b=o\), then every \(a\in A\) satisfies \(L_a\le R\).
\end{proposition}

\begin{proof}
Here \(R=\PP(U)\), whereas
\(L_a=\PP(o\con A,a\con o)\le\PP(U)\).
\end{proof}

\section{Zero score under positive avoidance}

The next lemma removes one boundary case from the open problem.  It is
stated without assuming that all edge weights lie strictly between zero
and one.

\begin{lemma}[Support-graph lemma]\label{lem:zero-score}
Assume \(\PP(D)>0\), let \(a\in A\), and suppose \(\rho_a=0\).  Then,
up to a null event,
\begin{equation}\label{eq:inclusion-zero}
 \{a\con b\}\subseteq\{o\con b,\ o\con A\}.
\end{equation}
In particular, \(L_a\le R\).
\end{lemma}

\begin{proof}
Let \(E_1=\{e:p_e=1\}\) be the forced-open edges and
\(E_*=\{e:0<p_e<1\}\) the genuinely random edges.  Every prescribed
subset of \(E_*\), together with all of \(E_1\), has positive probability.
Since \(\PP(D)>0\), the component of \(o\) in the forced-open graph
\((V,E_1)\) contains no point of \(A\).

Suppose that a supported configuration contains an open path from \(a\)
to \(b\), but does not connect \(o\) to \(b\).  Choose a simple open
\(a\)-to-\(b\) path \(P\).  Form a new supported configuration by retaining
all forced-open edges and retaining among the genuinely random edges only
those belonging to \(P\).  It still connects \(a\) to \(b\).

It cannot connect \(o\) to \(A\).  Indeed, a connection using no edge of
\(P\cap E_*\) would put a point of \(A\) in the forced-open component of
\(o\).  A connection using such an edge would connect \(o\) to the path
\(P\), and hence to \(b\), already in the original configuration.  Both
alternatives contradict the preceding paragraph or the choice of the
original configuration.  The new configuration therefore belongs to
\(\{a\con b\}\cap D\) and has positive probability, contradicting
\(\rho_a=0\).

Thus \(a\con b\) implies \(o\con b\) almost surely.  Whenever both
connections hold, \(o\con a\) through \(b\); since \(a\in A\), the event
\(U\) also holds.  This proves \eqref{eq:inclusion-zero}, and taking
probabilities gives \(L_a\le R\).
\end{proof}

\section{A pendant-relay reduction}

The strict-minimizer version and the positive-avoidance every-minimizer
version are essentially the same conjecture.  This reduces the remaining
work to a single precise target.

\begin{lemma}[Pendant tie breaker]\label{lem:pendant}
Assume \(\PP(D)>0\).  Let \(a\in A\) minimize \(\rho_x\), and suppose
\(\rho_a>0\).  For \(0<q<1\), add a new vertex \(\widehat a\), joined only
to \(a\) by an edge of probability \(q\), and replace \(a\) in the relay set:
\[
 \widehat A=(A\setminus\{a\})\cup\{\widehat a\}.
\]
Use hats for the quantities \eqref{eq:three-quantities} in the enlarged
graph.  Then \(\widehat a\) is the unique minimizer of the new Question 8
score.  More precisely,
\begin{align}
 \widehat\rho_{\widehat a}(q)&=q\rho_a,
 \label{eq:pendant-rhoa}\\
 \widehat\rho_x(q)&=q\rho_x+(1-q)
   \PP(x\con b,\ o\ncon A\setminus\{a\})
   \ge \rho_x \qquad(x\in A\setminus\{a\}),
 \label{eq:pendant-rhox}\\
 \widehat L_{\widehat a}(q)&=qL_a,
 \label{eq:pendant-L}\\
 \widehat R(q)&=qR+(1-q)
   \PP(o\con b,\ o\con A\setminus\{a\}).
 \label{eq:pendant-R}
\end{align}
\end{lemma}

\begin{proof}
Condition on the state of the new pendant edge.  If it is open, wiring
\(\widehat a\) to \(a\) identifies the new events with the original ones.
If it is closed, \(\widehat a\) cannot connect to the old target \(b\), and
the new observer-to-relay event is \(\{o\con A\setminus\{a\}\}\).
This proves all four identities.  Since
\[
 \{o\ncon A\}\subseteq\{o\ncon A\setminus\{a\}\},
\]
the last probability in \eqref{eq:pendant-rhox} is at least \(\rho_x\),
which proves the displayed lower bound.  Therefore, for every old relay
\(x\ne a\),
\[
 \widehat\rho_{\widehat a}=q\rho_a<\rho_a
 \le\rho_x\le\widehat\rho_x.
\]
Thus \(\widehat a\) is the unique minimizer.
\end{proof}

\begin{theorem}[Equivalence of the two open cores]\label{thm:equivalence}
The following global assertions are equivalent.
\begin{enumerate}[label=\textup{(\roman*)}]
\item The strict-minimizer form, Definition~\ref{def:strict}, holds for
every finite weighted graph.
\item The positive-avoidance form, Definition~\ref{def:positive}, holds for
every finite weighted graph and every choice among tied minimizers.
\end{enumerate}
\end{theorem}

\begin{proof}
Assume (i), suppose \(\PP(D)>0\), and choose any minimizing relay \(a\).
If \(\rho_a=0\), Lemma~\ref{lem:zero-score} proves \(L_a\le R\).  If
\(\rho_a>0\), apply (i) to the unique minimizer \(\widehat a\) constructed
in Lemma~\ref{lem:pendant}.  Equations \eqref{eq:pendant-L} and
\eqref{eq:pendant-R} give
\[
 qL_a\le qR+(1-q)\PP(o\con b,\ o\con A\setminus\{a\})
 \qquad(0<q<1).
\]
Letting \(q\uparrow1\) gives \(L_a\le R\).  No limit theorem is needed
for a formal proof: if \(\Delta=L_a-R>0\), then \(0<\Delta\le1\), and the
choice \(q=1-\Delta/4\) makes
\(q\Delta>(1-q)\PP(o\con b,o\con A\setminus\{a\})\), contradicting the
last display.  Hence (ii) holds.

Conversely, suppose (ii) and let \(a\) be a strict minimizer.  If
\(\PP(D)>0\), apply (ii).  If \(\PP(D)=0\), all the scores are zero, so a
strict minimizer can exist only when \(A=\{a\}\); Proposition~\ref{prop:singleton}
then applies.  Thus (i) holds.
\end{proof}

\begin{remark}
Theorem~\ref{thm:equivalence} does not restore the false every-minimizer
claim on \(\PP(D)=0\).  The positive-avoidance hypothesis is indispensable
in the stated equivalence.
\end{remark}

\section{Reduction to interior edge weights}

\begin{proposition}[Boundary closure for a strict minimizer]
\label{prop:interior}
It is enough to prove Definition~\ref{def:strict} when every edge weight
lies in \((0,1)\).
\end{proposition}

\begin{proof}
The claim is immediate for a singleton relay set.  Otherwise, a strict
minimizer has a positive gap
\[
 \delta:=\min_{x\in A\setminus\{a\}}(\rho_x-\rho_a)>0.
\]
Every event probability on a finite graph is a multilinear polynomial in
the edge weights.  Approximate an arbitrary weight vector \(p\in[0,1]^E\)
by vectors \(p^{(n)}\in(0,1)^E\).  By continuity, the same relay \(a\)
remains the strict minimizer for all sufficiently large \(n\).  The desired
inequality at \(p^{(n)}\), followed by another use of continuity, gives the
inequality at \(p\).
\end{proof}

Combining Proposition~\ref{prop:interior}, Theorem~\ref{thm:equivalence},
and Lemma~\ref{lem:zero-score}, the genuinely open case may be stated with
all of the following assumptions:
\begin{equation}\label{eq:open-core-assumptions}
 p_e\in(0,1),\quad |A|\ge2,\quad \PP(D)>0,\quad \rho_a>0,
 \quad \rho_a<\rho_x\ (x\ne a).
\end{equation}
The conclusion still needed is \(L_a\le R\).

\section{An exact event reduction and a proved sufficient condition}

Fix \(a\in A\), put \(T=A\setminus\{a\}\), and define
\begin{equation}\label{eq:Ea}
 E_a:=\{o\ncon a,\ o\con T\}.
\end{equation}

\begin{lemma}[Cancellation on the connected branch]\label{lem:cancel}
For every \(a\in A\),
\begin{equation}\label{eq:delta-core}
 L_a-R
 =\PP(a\con b,E_a)-\PP(o\con b,E_a).
\end{equation}
\end{lemma}

\begin{proof}
Partition \(U\) according as \(o\con a\) or \(o\ncon a\).  On
\(\{o\con a\}\), the clusters of \(o\) and \(a\) coincide, so the two
contributions cancel.  On \(\{o\ncon a\}\), the event \(U\) is exactly
\(\{o\con T\}\), giving \eqref{eq:delta-core}.
\end{proof}

The next proposition records a route already supported by the avoided
first-relay machinery.  It is a sufficient condition, not a proof that the
Question 8 selection rule implies the condition.

Let
\begin{equation}\label{eq:Tplus}
 T_+:=\{x\in T:\PP(x\ncon a)>0\}.
\end{equation}
For a relay \(x\in T_+\), set
\begin{equation}\label{eq:mx}
 m_x^{a}:=\PP(x\con b\mid x\ncon a).
\end{equation}
Choose an injective rank on \(T_+\) that is nondecreasing in these
conditional probabilities, breaking ties arbitrarily.  Let \(J_x\) be the
event that \(o\ncon a\) and \(x\) is the first relay of \(T_+\), in that
rank, reached by \(o\).  These events are disjoint and their union is
\(E_a\), up to a null event.  Indeed, if
\(x\in T\setminus T_+\), then
\(\{o\ncon a,o\con x\}\subseteq\{x\ncon a\}\) is null.  A finite union
over the inactive relays is still null, so deleting them makes no
contribution to \(E_a\).

\begin{proposition}[A formalizable sufficient criterion]
\label{prop:sufficient}
Assume \(\PP(o\ncon a)>0\).  If
\begin{equation}\label{eq:sufficient}
 \PP(a\con b,a\ncon o)\,\PP(E_a)
 \le \PP(a\ncon o)\sum_{x\in T_+}\PP(J_x)m_x^a,
\end{equation}
then \(L_a\le R\).  In particular, it is sufficient that
\begin{equation}\label{eq:strong-sufficient}
 \PP(a\con b\mid a\ncon o)
 \le \min_{x\in T_+}
       \PP(x\con b\mid x\ncon a),
\end{equation}
provided \(E_a\) has positive probability.
\end{proposition}

\begin{proof}
The assumption \(\PP(o\ncon a)>0\) in particular implies \(a\ne o\).
The denominator-free form of the two-cluster BHK inequality, conditioned
on \(o\ncon a\), gives
\begin{equation}\label{eq:BHK-use}
 \PP(o\ncon a)\PP(a\con b,E_a)
 \le \PP(a\con b,a\ncon o)\PP(E_a).
\end{equation}
The avoided first-relay inequality with avoided set \(\{a\}\), active
relay set \(T_+\), and cluster functional
\(K\mapsto\1_{\{b\in K\}}\) gives
\begin{equation}\label{eq:GEN-use}
 \PP(o\con b,E_a)\ge\sum_{x\in T_+}\PP(J_x)m_x^a.
\end{equation}
Combining \eqref{eq:sufficient}, \eqref{eq:BHK-use}, and
\eqref{eq:GEN-use}, then cancelling the positive factor
\(\PP(o\ncon a)\) and using Lemma~\ref{lem:cancel}, proves the first claim.
For \eqref{eq:strong-sufficient}, use
\(\sum_{x\in T_+}\PP(J_x)=\PP(E_a)\).  If \(E_a\) has positive
probability, then \(T_+\) is nonempty, so the minimum in
\eqref{eq:strong-sufficient} is defined.
\end{proof}

The missing bridge is now explicit: no proof is presently supplied that
the Question 8 hypothesis \(\rho_a\le\rho_x\) forces
\eqref{eq:sufficient}.  Numerical agreement of the final inequality cannot
replace this implication.

\subsection{A sharper conjectural bridge}

The following inequality is a useful research target, but it is
\emph{not proved in this note}.  Work under the strict assumptions
\eqref{eq:open-core-assumptions}.  For \(x\in T\), put
\begin{equation}\label{eq:beta}
 d_x:=\PP(x\ncon a),\qquad
 \beta_x:=\frac{\rho_x-\rho_a}{d_x}.
\end{equation}
Every \(d_x\) is positive: if \(x\con a\) almost surely, then the
indicators of \(x\con b\) and \(a\con b\) agree almost surely, forcing
\(\rho_x=\rho_a\) and contradicting strictness.  Rank \(T\) increasingly
by \(\beta_x\), and define the associated first-relay events \(J_x\) on
\(\{o\ncon a\}\) as above.  The candidate inequality is
\begin{equation}\label{eq:conjectural-bridge}
 \boxed{
 R-L_a\ \ge\ \sum_{x\in T}\PP(J_x)
       \frac{\rho_x-\rho_a}{\PP(x\ncon a)}.}
\end{equation}
It would immediately settle the strict case because every term on the
right is nonnegative.  It has survived small-graph computational tests,
but that is evidence only.

The direct attempted derivation has a specific obstruction.  The natural
cluster projection whose avoided relay integral equals
\(\rho_x-\rho_a\) is not monotone, so the existing avoided first-relay
theorem cannot be applied to it.  Any proof of
\eqref{eq:conjectural-bridge} therefore needs either a monotone correction
whose error has a controlled sign or a new signed version of the
first-relay inequality.  A Lean development may state
\eqref{eq:conjectural-bridge} as a closed proposition for future work, but
must not import it as an axiom or use it in the proved chain above.

\section{Lean handoff}

The formalization should keep false, proved, and open declarations in
separate namespaces or files.  A minimal declaration list is as follows.

\begin{enumerate}[label=\texttt{Q8.\arabic*},leftmargin=4em]
\item \texttt{U}, \texttt{q8L}, \texttt{q8Score}, and \texttt{q8R},
corresponding to \eqref{eq:U}--\eqref{eq:three-quantities}.
\item \texttt{IsQ8Min} and \texttt{IsQ8StrictMin}.
\item Closed propositions \texttt{Question8EveryMin},
\texttt{Question8Strict}, \texttt{Question8Positive}, and
\texttt{Question8ExistsGoodMin}.  The first is to be refuted; the other
three are statements only.
\item \texttt{q8\_total\_split}, the identity \eqref{eq:total-split}.
\item The six identities \eqref{eq:path1}--\eqref{eq:path3} on
\texttt{Fin 4}, followed by \texttt{not\_question8EveryMin}.
\item \texttt{q8\_singleton}, \texttt{q8\_bEqO}, and
\texttt{q8\_degenerate\_allMin}.
\item A support-atom lemma: a configuration has positive product-Bernoulli
mass exactly when it contains every pair of weight one and contains no pair
of weight zero.  Use it in \texttt{q8\_zeroScore},
Lemma~\ref{lem:zero-score}, rather than introducing conditional
probabilities.
\item The four pendant identities \eqref{eq:pendant-rhoa}--\eqref{eq:pendant-R},
then \texttt{q8Strict\_iff\_q8Positive}.  Concretely, enlarge
\texttt{Fin n} to \texttt{Fin (n+1)}, embed old vertices with
\texttt{Fin.castSucc}, use \texttt{Fin.last n} for \(\widehat a\), and
define the new \texttt{Sym2} weight by cases: the pendant pair has weight
\(q\), embedded old pairs retain \(w\), and all other new pairs have weight
zero.
\item \texttt{q8\_target\_split}, identity \eqref{eq:delta-core}.
\item \texttt{q8\_of\_firstRelayCriterion},
Proposition~\ref{prop:sufficient}, importing the already formalized avoided
first-relay theorems \texttt{KNAll.genY\_all} and
\texttt{sum\_le\_setIntegral\_of\_genY}, together with
\texttt{BHK2006\_twoClusterConditionalAssociation.negCorrelation}.
\item A statement-only declaration \texttt{Question8Bridge} for
\eqref{eq:conjectural-bridge}.  It is not a theorem and must have no role in
the proofs of the preceding declarations.
\end{enumerate}

For Proposition~\ref{prop:interior}, reuse
\texttt{prodBernoulli\_real\_continuous}; the positive strict-minimum gap
keeps the same relay minimizing along the approximating sequence.

No theorem named \texttt{question8\_holds} should be introduced unless the
open core \eqref{eq:open-core-assumptions} is actually proved.  In
particular, the theorem
\[
 R\ge\min_{x\in A}L_x
\]
must not be used with a \(\rho\)-minimizer without an additional argument
relating the two rankings.

\subsection*{Suggested dependency graph}

\[
\begin{array}{c}
\text{event definitions}
\\[2mm]
\swarrow\qquad\downarrow\qquad\searrow
\\[-1mm]
\text{path identities}
\qquad
\text{support lemma}
\qquad
\text{pendant identities}
\\[1mm]
\downarrow\hspace{35mm}\searrow\qquad\swarrow
\\[-1mm]
\neg\text{EveryMin}
\hspace{24mm}
\text{Strict}\Longleftrightarrow\text{PositiveAvoidance}
\end{array}
\]

Separately,
\[
\text{BHK}+\text{avoided first-relay}+\eqref{eq:delta-core}
\Longrightarrow\text{Proposition~\ref{prop:sufficient}}.
\]
This organization lets a checker certify every result in the present note
without inserting a placeholder proof for the genuinely open implication.

\section{Final status}

The exact conclusion is:
\begin{enumerate}[label=\arabic*.]
\item The printed score is \(\rho_x=\PP(x\con b,o\ncon A)\).
\item The every-minimizer reading over weights in \([0,1]\) is false by an
exact four-vertex counterexample.
\item The previous manuscript proved the different statement obtained by
minimizing \(L_x=\PP(o\con A,x\con b)\).
\item The strict-minimizer version and the every-minimizer version under
\(\PP(o\ncon A)>0\) are equivalent, but no complete proof of either is
given here.
\item The remaining proof obligation is the implication from the strict
ordering of the \(\rho_x\)'s in \eqref{eq:open-core-assumptions} to
\(L_a\le R\), equivalently the inequality in
Lemma~\ref{lem:cancel}.
\end{enumerate}

\begin{thebibliography}{1}
\bibitem{KN}
G.~Kozma and S.~Nitzan,
\emph{A reduction of the $\theta(p_c)=0$ problem to a conjectured
inequality}, arXiv:2401.12397, 2024.  Question 8 appears on p.~36.
\end{thebibliography}

\end{document}
